\begin{lemma}[Replaying an augmentation sequence through an ordered union]
Let $P,L,R,U$ be tree-sequences for the same finite loopless multigraph
and distinguished first edge, and let $\varphi$ be a fixed partial edge
coloring.  Suppose that $R$ is obtained from $P$ by a
$\varphi$-\noderef{TashkinovAugmentationSequence}, that $P$ is an exact
ordered \noderef{TreeSequenceInitialSegment} of $L$, and that $U$ is
an \noderef{OrderedTreeUnion} of $L$ and $R$.  Then $U$ is obtained
from $L$ by a $\varphi$-Tashkinov augmentation sequence.
\end{lemma}
\begin{proof}
Put $\ell=|L|$.  For every $m$ with
$\ell\leq m\leq |U|$, let $U_m$ be the exact ordered prefix of $U$
having $m$ edge entries, as in \noderef{TreeSequencePrefix}.  We prove
inductively that $U_m$ is obtained from $L$ by a
$\varphi$-\noderef{TashkinovAugmentationSequence}.

The left-segment clause of \noderef{OrderedTreeUnion} says that the first
$\ell$ edge entries and first $\ell+1$ vertex entries of $U$ are
exactly those of $L$.  Hence $U_\ell=L$ as ordered edge-valued
tree-sequences, and the empty augmentation sequence proves the base case.

Fix $\ell\leq m<|U|$ and assume the assertion for $U_m$.  Let $j=m$
be the index of the next edge entry of $U$.  The
``new entry from the right'' clause of \noderef{OrderedTreeUnion} supplies
an index $i<|R|$ such that
\[
 U.e_j=R.e_i,
 \qquad U.y_{j+1}=R.y_{i+1}.                                  \tag{1}
\]
The right vertex in (1) is not in $V(L)$.  Indeed, if it were the
$a$-th vertex of $L$, the left-segment equality would make it the
$a$-th vertex of $U$.  Equation (1) would then identify the distinct
vertex indices $j+1>\ell\geq a$, contrary to the vertex injectivity in
\noderef{TreeSequence}.

By \noderef{TashkinovAugmentationSequenceInitialSegment}, $P$ is an exact
ordered initial segment of $R$.  Since it is also an exact ordered
initial segment of $L$, the preceding exclusion implies
$i\geq |P|$: if $i<|P|$, then $R.y_{i+1}=P.y_{i+1}$ would already
belong to $L$.

Let $R_i$ be the prefix of $R$ having $i$ edge entries.  The
inequality just proved makes $P$ an exact prefix of $R_i$, while
$R_i$ is an exact prefix of $R$.  Apply
\noderef{TashkinovAugmentationSequenceSuffix} to the actual suffix
$R_i\preccurlyeq R$.  It gives a $\varphi$-augmentation sequence from
$R_i$ to $R$.  This suffix is nonempty because $i<|R|$, so isolate
its first one-step augmentation $R_i\to C$.  By
\noderef{TashkinovAugmentation}, $C$ has length $i+1$, its new edge
has some color $c$, and
\[
                    c\in\overline\varphi(R_i).                 \tag{2}
\]
The remaining suffix from $C$ to $R$, together with
\noderef{TashkinovAugmentationSequenceInitialSegment}, makes $C$ an
exact prefix of $R$.  Therefore the exact new edge in this first step
is $R.e_i$.  In particular,
\[
                         \varphi(R.e_i)=c.                      \tag{3}
\]

We next prove that $c$ is missing on $U_m$.  By (2) and
\noderef{TreeMissingColors}, choose a vertex $v\in V(R_i)$ at which
$c$ is missing.  If $v\in V(L)$, the left-segment equality places
$v$ among the first $\ell+1\leq m+1$ vertices of $U$, hence in
$V(U_m)$.

Suppose instead that $v\notin V(L)$.  The initial vertex of $R_i$ is
the initial vertex of $P$, hence is in $L$; thus
$v=R.y_{r+1}$ for some $r<i$.  The ``every new right entry'' clause of
\noderef{OrderedTreeUnion} supplies an index $q\geq\ell$ at which the
exact right entry $(R.e_r,R.y_{r+1})$ occurs in $U$.  Both
$R.y_{r+1}$ and $R.y_{i+1}$ are outside $L$, and $r<i$, so the
``right order'' clause gives $q<j=m$.  Hence $v=U.y_{q+1}$ belongs to
$U_m$.  In either case the same vertex $v$ witnesses
\[
                         c\in\overline\varphi(U_m).             \tag{4}
\]

The prefix $U_m$ is an exact ordered initial segment of $U_{m+1}$,
and their lengths differ by one.  Equations (1) and (3) identify the new
edge of $U_{m+1}$ as the same exact edge identity $R.e_i$ of color
$c$; equation (4) says that $c$ is missing on $U_m$.  Thus all
clauses of \noderef{TashkinovAugmentation} hold for
$U_m\to U_{m+1}$.  Appending this step to the induction hypothesis
proves the assertion for $m+1$.

Taking $m=|U|$, the prefix $U_m$ is the full ordered tree-sequence
$U$.  The resulting chain is the required
$\varphi$-Tashkinov augmentation sequence from $L$ to $U$.
\end{proof}
